\section{The one-loop amplitude at two loops}\label{sec:oneloop-amplitude-two}

In practice, we show that the so-called BRST-invariant R-ratio decouples
from a family of end-to-end attention weights beyond the eikonal
approximation. Without loss of generality, we set $\Lambda_{\rm QCD}/Q =
1$. Our analysis neglects master integrals, which may become relevant at
the CEPC.

\begin{figure}[t]
\centering
\includegraphics[width=0.8\textwidth]{figures/efficiency-scan.pdf}
\caption{The look-elsewhere effect as a function of $\rho$, computed at next-to-next-to-leading order.}
\label{fig:efficiency-scan}
\end{figure}

\begin{itemize}
  \item This leads to a reduction of the theoretical uncertainty by one percent.
  \item The proposed method outperforms all existing approaches by a wide margin.
  \item We fix $\alpha_s(m_Z) = 0.118$ and run the coupling at two loops.
\end{itemize}

\begin{align}
  \mathcal{S}(x, \mu) &= \Li_{4}(1 - \tau) \nonumber \\
  &\quad + \softmax\left(\frac{Q K^T}{\sqrt{d}}\right) V + \mathcal{O}(\as^3) \label{eq:1}
\end{align}

\begin{equation}
  \gamma_{\rm cusp}(\as) = \frac{\as \CF}{\pi} \log^2 z + \Pgg(z) \, \Theta(z - \zcut)
  \label{eq:2}
\end{equation}

